\section{Matter and Subtle Bodies}

\begin{quotex}
The purely biological or scientific standpoint falls short in psychology because it is intellectual only. But the psychic phenomenon cannot be grasped in its totality by the intellect, for it consists not only of meaning but also of value, and this depends on the intensity of the accompanying feeling-tones. The affective value gives the measure of the intensity of an idea, and the intensity in its turn expresses that idea's energic tension, its effective potential. \flright{\textsc{Carl Jung}, \emph{Aion}}

\end{quotex}
The common illusion of many is the claim that they respect ``facts". Facts are value-neutral because they are grasped by the intellect which presumably judges impartially. However, facts involving humans elicit a meaning, to which a value is attached. Then a positive valuation is pursued and a negative one is eschewed. Hence, the valuation is accompanied with a more or less strong emotion.

Few are able to separate the intellectual content of a fact from its meaning and valuation. The result is usually a strong negative emotion like anger, fear, worry, depression, etc., which is falsely attributed to the ``fact". They fail to see that the meaning they believe they see out there is entirely subjective. So they fail to understand the real reason for their upset. Often it is because they only see the past thereby missing the experience of the present.

To sum up: \emph{Thought} leads to \emph{Feelings} leads to \emph{Will}. If your thoughts are incorrect, you will feel bad and do the wrong thing.

\paragraph{Naturalism}
Naturalism is the theory that all existence can be explained in terms of the motion of particles. These theories go at least as far back as the Pre-Socratics and is still popular today among the educated. For example, see the physicist Sabine Hossenfelder\footnote{\url{https://www.youtube.com/watch?v=zpU_e3jh_FY}} discus free will. She adds the wrinkle of differential equations which presumes that reality is continuous.

Motion requires time and space. So naturalism admits three preconditions for manifestation—space, time and matter—but denies two others: form and psyche (or life, consciousness, intellect).

\paragraph{One}
Plato fixed things by understanding the missing two. There are no things without a form to distinguish them. The accidental arrangements of matter or particles cannot create a new form. Rather, matter participates in a form. Hence he was led to believe in a realm of Forms or Ideas. The highest of the Ideas is the idea of the One, which is also the greatest good.

Things have more being, or power, as they participate in the One. That is, the oneness is essential to the thing. Natural kinds like a tree, a squirrel or a human being have a natural unity. However, humans are internally divided, being inconsistent from day to day; they make and break promises, change their minds, are full of doubts, and so on. Said differently, they do not possess a true Self which would give unity to their being. As a person works to become more unified in consciousness, he becomes more like the One, i.e., he has ``more being".

Artifacts, on the other hand, do not have an inner unity. Rather, their unity is imposed ``from the outside", i.e., by a human designer. Automobiles, can openers, toilets, and androids are artifacts.

\paragraph{Subtle Matter}
Traditional teachings recognize several levels of being, from prime matter to more subtle levels, such as the astral or ethereal levels. Wolfgang Smith recognizes that physicists have discovered a level above prime matter, populated with subatomic and atomic particles, but below the corporeal level of macroscopic objects; they are intimately related, yet independent. On the one hand, it is certainly possible to conceive of a physical world devoid of any actual corporeal things, and on the other to conceive of corporeal objects apart from the physical level. For example, the ``resurrection body", for all appearances, looks like a physical body, yet does not follow the laws of physics.

Such subtle realms used to be part of common knowledge, but not so much today. ``Only matter is real," they say. But now what matter ``is" is not so clear. Matter is equivalent to energy, being condensed energy, yet there is no non-circular definition of energy.

If we turn to quantum mechanics, then according to Schrödinger's equation, matter is ultimately in a state of indeterminacy; that is, it is nothing in particular. But that is what the ancients believed, even without the famous equation. How it becomes specific macroscopic things is rather unclear (the so-called measurement problem).

General Relativity does not help much, since it denies any force called ``gravity". Rather, massive objects distort the geometry of the universe by creating holes in a fourth spatial dimension. No one can see that dimension, so that a straight line path in that geometry falsely appears to be attracted to such objects. So, what is the point of mass, if masses don't exert forces on each other? The universe is just lumpy, and we experience those lumps as masses.

\paragraph{The Planetary Spheres}
In our vanity, we presume that the ancients mistakenly believed that the earth was surrounded by the planetary spheres. But is that false or are we blinded by science? Actually, physics doesn't suffer at all if we assume that the planets revolve around the Sun or the Earth, although it could get mathematically messy. Einstein himself conceded that no theory or observation could resolve the question.

Nevertheless, Dante was one such ancient who actually described his journey though the planetary spheres. He certainly was not describing a space journey in a rocket ship or flying saucers, although it involved a flight. Dante flies in his ethereal body, which is naturally attracted to God when freed from the physical body. In a sense, then, Dante was having an ``out of body experience", although he experienced much more and in greater depth, than what is often described today in such experiences.

So we can speculate that the prime elements are not merely symbols for the physical states of matter, but actually represent different states of being:

\begin{itemize}
\item \textbf{Earth}: physical body 
\item \textbf{Water}: life body 
\item \textbf{Fire}: astral body 
\item \textbf{Air}: intellectual soul 
\item \textbf{Ether}: spiritual body 
\end{itemize}


\flrightit{Posted on 2020-11-05 by Cologero }
